\section*{Abstract}

University teaching in Bangladesh is usually delivered in a mixture of Bengali and
English, and students write that mixture in the Roman alphabet, a form commonly
called Banglish. Existing speech recognition does not serve it. Given a recording of
such a lecture, a multilingual model produces an English translation of what was
said, and a Bengali model writes the Bengali script that no student in the class
uses. The whiteboard, which carries most of the worked examples in a class taught
without slides, is lost entirely.

This thesis presents InsightLens, a system that takes a recorded Banglish whiteboard
lecture and produces a lecture note containing both the speech and the board. It
works in four stages. A Whisper large-v3-turbo model adapted with Low-Rank Adaptation
writes a timestamped transcript in romanized Banglish. Each whiteboard is
reconstructed from the video by tiling the frame and taking every tile from a moment
when the lecturer was not in front of it, so that every pixel of the result is
unmodified camera output. Numbered coloured boxes are drawn around each block of
writing and a vision-language model names and transcribes each box. A language model
then writes one section of notes per board, pointing the reader at the boxes, with
every quotation checked word for word against the transcript.

A corpus of 44 lecture recordings from three lecturers was collected, of which 28
lectures covering 5.15 hours were transcribed by hand against a written transcription
standard. On six whole lectures held out at random, fine-tuning cuts the median
per-clip character error rate from 67.7 to 15.8 and 16.0 per cent in two random seeds
and the word error rate from 93.9 to 42.5 and 41.7 per cent, better on 164 and 167 of
177 clips. Holding out an entire lecturer, the harder design, gives 72.8 to 50.3 per
cent character error in all six runs. Against hand-verified answer keys covering 45
boards and 436 items, changing the board-reading prompt from keywords to full
transcription raises recall from 31.2 to 88.8 per cent, while changing the model from
3 to 7 billion parameters is worth 5.3 points, so the prompt matters far more than the
model size. The generated notes carry 89.1 per cent of the board content against 37.2
per cent for a conventional summarisation prompt.

Results that did not support the design are reported in the same detail. Dual
recogniser fusion, biasing the decoder towards board words and using occlusion as a
pointer all failed. One training seed at the pre-registered learning rate diverged
completely. Measures assembled out of English strings were found, twice and
independently, to reward a model that translates code-mixed speech and to penalise one
that transcribes it, which is offered as a methodological finding for anyone
evaluating systems of this kind. The notes have not been evaluated with readers, and
that is the first item of future work.

\vspace{1cm}
\textbf{Keywords:} Code-mixed speech recognition; Banglish; Bengali-English;
Low-Rank Adaptation; Whisper; vision-language models; whiteboard reconstruction;
lecture note generation; Set-of-Mark prompting; low-resource speech
\pagebreak
